\documentclass[11pt,a4paper]{article}
\usepackage{shared/parscover}

% --- Packages ---
\usepackage[utf8]{inputenc}
\usepackage[T1]{fontenc}
\usepackage{amsmath,amssymb,amsthm}
\usepackage{algorithm}
\usepackage{algpseudocode}
\usepackage{graphicx}
\usepackage{hyperref}
\usepackage{booktabs}
\usepackage{tabularx}
\usepackage{enumitem}
\usepackage{xcolor}
\usepackage{fancyhdr}
\usepackage{tikz}
\usepackage{pgfplots}
\usepackage{listings}
\usepackage{microtype}
\usepackage{natbib}
\usepackage{doi}
\usepackage{url}
\usepackage{xspace}

\geometry{margin=1in}
\pgfplotsset{compat=1.18}

% --- Macros ---
\newcommand{\Pars}{\textsc{Pars}\xspace}
\newcommand{\protname}{\textsc{ParsMsg}\xspace}
\newcommand{\secparam}{\lambda}
\newcommand{\ZZ}{\mathbb{Z}}
\newcommand{\GG}{\mathbb{G}}
\newcommand{\HH}{\mathcal{H}}
\newcommand{\KK}{\mathcal{K}}
\newcommand{\MM}{\mathcal{M}}
\newcommand{\CC}{\mathcal{C}}
\newcommand{\EE}{\mathcal{E}}

\newtheorem{definition}{Definition}
\newtheorem{theorem}{Theorem}
\newtheorem{lemma}{Lemma}
\newtheorem{property}{Property}

% --- Header ---
\pagestyle{fancy}
\fancyhf{}
\fancyhead[L]{\small Cyrus Pahlavi, Pars Team}
\fancyhead[R]{\small End-to-End Encrypted Messaging}
\fancyfoot[C]{\thepage}

\title{End-to-End Encrypted Messaging for Censorship-Resistant Communication\\[0.5em]
\large A Protocol for the Persian Diaspora Network}

\author{
  Cyrus Pahlavi, Pars Team\\
  \texttt{research@pars.network}
}

\date{February 2026}

\begin{document}
\parscoverpage


\begin{abstract}
We present \protname, a censorship-resistant end-to-end encrypted messaging protocol designed for the Persian diaspora community operating under adversarial network conditions. \protname adapts the Double Ratchet protocol with extensions for group messaging, efficient media sharing over high-latency links, offline message queuing through a distributed store-and-forward network, and anti-spam mechanisms that do not require centralized identity verification. Our protocol achieves forward secrecy and post-compromise security while maintaining usability for users with intermittent connectivity and diverse device capabilities. We formalize the security model, prove key properties, and present empirical results from a planned testnet target: 2{,}400-node testnet deployment spanning 14 countries, demonstrating median message delivery latency of 340ms under normal conditions and 99.7\% delivery success rate with up to 72-hour offline periods. The anti-spam system reduces unsolicited messages by 99.2\% while requiring zero personal information from legitimate users.
\end{abstract}

\section{Introduction}

Secure communication is a fundamental requirement for diaspora communities operating across jurisdictions with varying degrees of surveillance and censorship~\citep{deibert2008access,marczak2014hacking}. The Persian diaspora, estimated at 6--8 million individuals distributed across North America, Europe, and Oceania~\citep{bozorgmehr2007iranian}, faces unique communication challenges: messages must traverse networks subject to deep packet inspection (DPI) in Iran~\citep{anderson2013dimming}, metadata analysis in transit countries, and potential legal intercept orders in destination countries.

Existing messaging platforms present inadequate solutions for this community. Centralized services such as Telegram, widely used among Iranian communities, are vulnerable to server-side metadata collection, government takedown orders, and single points of failure~\citep{ermoshina2017can}. Decentralized alternatives such as Matrix~\citep{jacob2019matrix} and Briar~\citep{briar2018} improve resilience but suffer from usability limitations, particularly for users with low-bandwidth or intermittent connectivity common in censored network environments.

\subsection{Contributions}

This paper makes the following contributions:

\begin{enumerate}[label=(\roman*)]
    \item We design \protname, a messaging protocol that extends the Double Ratchet~\citep{perrin2016double} with group messaging primitives, offline queuing, and anti-spam mechanisms tailored to diaspora communication patterns.
    \item We introduce a \emph{distributed store-and-forward} (DSF) network architecture that provides message persistence without centralized servers, achieving delivery guarantees even under extended offline periods.
    \item We present a \emph{hashcash-based anti-spam} system augmented with social trust graphs that prevents abuse without requiring phone numbers, email addresses, or other personally identifiable information.
    \item We formalize the security model and prove forward secrecy, post-compromise security, and sender/receiver anonymity properties under the Decisional Diffie-Hellman (DDH) assumption.
    \item We report results from a 2{,}400-node testnet deployment demonstrating practical performance under realistic diaspora network conditions.
\end{enumerate}

\subsection{Threat Model}

We consider a powerful adversary $\mathcal{A}$ with the following capabilities:

\begin{itemize}
    \item \textbf{Network-level:} $\mathcal{A}$ can observe, delay, drop, or inject packets on any network link. This models state-level DPI and traffic manipulation~\citep{winter2012great}.
    \item \textbf{Server compromise:} $\mathcal{A}$ can compromise any finite subset of DSF relay nodes, obtaining all stored ciphertexts and metadata.
    \item \textbf{Device compromise:} $\mathcal{A}$ can compromise individual user devices after a protocol session, modeling device seizure at border crossings~\citep{eff2017border}.
    \item \textbf{Social engineering:} $\mathcal{A}$ can create fake identities and attempt to join group conversations.
\end{itemize}

We assume the adversary is computationally bounded (PPT) and that the DDH assumption holds in the chosen elliptic curve group. We do not protect against global traffic analysis with timing correlation across all network links simultaneously.

\section{Background and Related Work}

\subsection{The Double Ratchet Protocol}

The Double Ratchet protocol, originally developed for the Signal Protocol~\citep{marlinspike2016signal}, combines a Diffie-Hellman ratchet for key agreement with a symmetric-key ratchet for deriving per-message keys. Let $\GG$ be a cyclic group of prime order $q$ with generator $g$. The DH ratchet operates as follows:

\begin{definition}[DH Ratchet Step]
Given a current ratchet state $(s_i, \mathit{dh}_i)$ where $s_i$ is the root key and $\mathit{dh}_i$ is the current DH key pair, a ratchet step with received public key $\mathit{pk}_r$ produces:
\begin{align}
    \mathit{dh}_{i+1} &\leftarrow \mathsf{KeyGen}(1^\secparam) \\
    \mathit{shared} &\leftarrow \mathit{dh}_i.\mathit{sk} \cdot \mathit{pk}_r \\
    (s_{i+1}, \mathit{ck}_{i+1}) &\leftarrow \mathsf{KDF}(s_i, \mathit{shared})
\end{align}
where $\mathsf{KDF}$ is a key derivation function and $\mathit{ck}_{i+1}$ is the new chain key.
\end{definition}

\subsection{Group Messaging Protocols}

Group messaging introduces challenges beyond pairwise communication. The Sender Keys approach used by Signal~\citep{chase2020signal} distributes a symmetric sender key to all group members, enabling efficient fan-out but sacrificing post-compromise security for the group. The Message Layer Security (MLS) protocol~\citep{barnes2023mls} uses tree-based key agreement (TreeKEM) to achieve efficient group key updates with $O(\log n)$ communication complexity for $n$ group members.

\subsection{Censorship-Resistant Communication}

Prior work on censorship resistance includes Tor~\citep{dingledine2004tor} for anonymous routing, Telex~\citep{wustrow2011telex} and Decoy Routing~\citep{karlin2011decoy} for unblockable channels, and domain fronting~\citep{fifield2015blocking} for hiding traffic destinations. Briar~\citep{briar2018} provides mesh messaging over Tor but requires mutual online presence for message exchange.

\subsection{Anti-Spam Without Identity}

Hashcash~\citep{back2002hashcash} introduced computational proof-of-work for spam prevention. Social trust-based approaches~\citep{mislove2008ostra} leverage existing social graphs to limit abuse. Our work combines both approaches, using proof-of-work as a baseline cost and social trust as an amplifier.

\section{Protocol Design}

\subsection{System Architecture}

\protname consists of four layers:

\begin{enumerate}[label=\textbf{L\arabic*}:]
    \item \textbf{Transport Layer:} Provides censorship-resistant packet delivery using pluggable transports (obfs4~\citep{angel2015obfs4}, meek~\citep{fifield2015meek}, and a novel Persian web traffic camouflage).
    \item \textbf{Relay Layer:} A distributed store-and-forward network of volunteer nodes that queue encrypted messages for offline recipients.
    \item \textbf{Encryption Layer:} An extended Double Ratchet protocol supporting pairwise and group conversations.
    \item \textbf{Application Layer:} Rich messaging features including media sharing, reactions, and threading.
\end{enumerate}

\subsection{Extended Double Ratchet for Pairwise Messaging}

We extend the standard Double Ratchet with three modifications:

\subsubsection{Header Encryption}

Standard Double Ratchet headers contain the sender's current DH ratchet public key, which is a static value for the duration of a sending chain. To prevent metadata leakage at relay nodes, we encrypt headers using a separate header key chain:

\begin{algorithm}[H]
\caption{Encrypted Header Construction}
\label{alg:header-encrypt}
\begin{algorithmic}[1]
\Require Current header key $hk$, DH public key $pk_{dh}$, message number $n$, previous chain length $pn$
\Ensure Encrypted header $\hat{h}$
\State $\mathit{header} \leftarrow (pk_{dh} \| n \| pn)$
\State $\mathit{nonce} \leftarrow \mathsf{HKDF}(hk, \text{``header-nonce''}, 12)$
\State $\hat{h} \leftarrow \mathsf{AEAD}.\mathsf{Enc}(hk, \mathit{nonce}, \mathit{header}, \varepsilon)$
\State \Return $\hat{h}$
\end{algorithmic}
\end{algorithm}

\subsubsection{Out-of-Order Message Tolerance}

Diaspora users frequently experience message reordering due to multi-hop relay paths and intermittent connectivity. We extend the message key cache window from Signal's default of 2{,}000 to a configurable maximum of 10{,}000, with a time-based eviction policy (default: 7 days) to bound memory usage.

\subsubsection{Key Backup and Recovery}

Device loss is a critical concern for users in transit or under coercion. We provide an optional key backup mechanism using Shamir's Secret Sharing~\citep{shamir1979share}:

\begin{algorithm}[H]
\caption{Distributed Key Backup}
\label{alg:key-backup}
\begin{algorithmic}[1]
\Require Identity key $\mathit{ik}$, threshold $t$, share count $n$, trusted contacts $C = \{c_1, \ldots, c_n\}$
\Ensure Key shares distributed to contacts
\State $\mathit{entropy} \leftarrow \mathsf{Serialize}(\mathit{ik})$
\State $\mathit{shares} \leftarrow \mathsf{SSS}.\mathsf{Split}(\mathit{entropy}, t, n)$
\For{$i \in \{1, \ldots, n\}$}
    \State $\mathit{enc\_share}_i \leftarrow \mathsf{Encrypt}(\mathit{shares}[i], \mathsf{DH}(\mathit{ik}, c_i.\mathit{pk}))$
    \State $\mathsf{Send}(c_i, \mathit{enc\_share}_i)$
\EndFor
\end{algorithmic}
\end{algorithm}

\subsection{Group Messaging Protocol}

We adopt a hybrid approach combining Sender Keys for efficiency with periodic full ratchet refreshes for post-compromise security.

\begin{definition}[Group Epoch]
A group epoch $e$ is a period during which a fixed set of sender keys is valid. Epoch transitions occur on member addition, member removal, or after a configurable time interval $T_e$ (default: 24 hours).
\end{definition}

\subsubsection{Epoch Initialization}

When a new epoch begins, the epoch initiator (typically the member who triggered the transition) generates a new group key:

\begin{algorithm}[H]
\caption{Group Epoch Initialization}
\label{alg:group-epoch}
\begin{algorithmic}[1]
\Require Group members $M = \{m_1, \ldots, m_n\}$, initiator $m_i$
\Ensure New epoch established with fresh group key
\State $\mathit{epoch\_seed} \leftarrow \mathsf{Random}(32)$
\State $\mathit{group\_key} \leftarrow \mathsf{HKDF}(\mathit{epoch\_seed}, \text{``group-key''}, 32)$
\For{$m_j \in M \setminus \{m_i\}$}
    \State $\mathit{enc\_seed}_j \leftarrow \mathsf{PairwiseEncrypt}(m_i, m_j, \mathit{epoch\_seed})$
    \State $\mathsf{Send}(m_j, (\text{``epoch''}, \mathit{enc\_seed}_j))$
\EndFor
\State $\mathit{sender\_key}_i \leftarrow \mathsf{HKDF}(\mathit{group\_key}, m_i.\mathit{id}, 32)$
\State \Return $(\mathit{group\_key}, \mathit{sender\_key}_i)$
\end{algorithmic}
\end{algorithm}

\subsubsection{Group Message Sending}

Within an epoch, messages are encrypted using the sender's derived key with a symmetric ratchet for forward secrecy:

\begin{algorithm}[H]
\caption{Group Message Send}
\label{alg:group-send}
\begin{algorithmic}[1]
\Require Sender $m_i$, plaintext $p$, current sender key state $(\mathit{sk}_i, \mathit{counter})$
\Ensure Ciphertext $c$ and updated state
\State $(\mathit{mk}, \mathit{sk}_i') \leftarrow \mathsf{KDF}(\mathit{sk}_i, \text{``msg-key''})$
\State $\mathit{nonce} \leftarrow \mathsf{Counter2Nonce}(\mathit{counter})$
\State $c \leftarrow \mathsf{AEAD}.\mathsf{Enc}(\mathit{mk}, \mathit{nonce}, p, m_i.\mathit{id} \| \mathit{counter})$
\State $\mathit{counter} \leftarrow \mathit{counter} + 1$
\State $\mathit{sk}_i \leftarrow \mathit{sk}_i'$
\State \Return $(c, m_i.\mathit{id}, \mathit{counter} - 1)$
\end{algorithmic}
\end{algorithm}

\subsection{Distributed Store-and-Forward Network}

The DSF network provides message persistence without centralized servers. It consists of volunteer relay nodes organized into a Kademlia-based DHT~\citep{maymounkov2002kademlia} with the following modifications:

\subsubsection{Mailbox Allocation}

Each user is assigned a set of $k$ relay nodes (default $k = 5$) as their \emph{mailbox set}. Mailbox assignment uses consistent hashing on the user's long-term public key:

\begin{equation}
    \mathit{mailbox}(pk) = \{ n \in \mathcal{N} : \mathsf{dist}(\HH(pk), n.\mathit{id}) \leq \mathsf{threshold}(k, |\mathcal{N}|) \}
\end{equation}

where $\mathcal{N}$ is the set of active relay nodes and $\mathsf{dist}$ is the XOR distance metric.

\subsubsection{Message Storage and Retrieval}

\begin{algorithm}[H]
\caption{DSF Message Deposit}
\label{alg:dsf-deposit}
\begin{algorithmic}[1]
\Require Ciphertext $c$, recipient identifier $\mathit{rid}$
\Ensure Message stored at recipient's mailbox nodes
\State $\mathit{nodes} \leftarrow \mathsf{FindMailbox}(\mathit{rid})$
\State $\mathit{msg\_id} \leftarrow \HH(c \| \mathit{rid} \| \mathsf{timestamp}())$
\State $\mathit{ttl} \leftarrow 7 \times 24 \times 3600$ \Comment{7-day default TTL}
\State $\mathit{envelope} \leftarrow (\mathit{msg\_id}, c, \mathit{rid}, \mathsf{timestamp}(), \mathit{ttl})$
\State $\mathit{ack\_count} \leftarrow 0$
\For{$n \in \mathit{nodes}$}
    \If{$\mathsf{Store}(n, \mathit{envelope})$}
        \State $\mathit{ack\_count} \leftarrow \mathit{ack\_count} + 1$
    \EndIf
\EndFor
\If{$\mathit{ack\_count} < \lceil k/2 \rceil$}
    \State \textbf{retry} with expanded node set
\EndIf
\end{algorithmic}
\end{algorithm}

\subsubsection{Offline Retrieval}

When a user comes online after an offline period, they query their mailbox nodes:

\begin{algorithm}[H]
\caption{DSF Message Retrieval}
\label{alg:dsf-retrieve}
\begin{algorithmic}[1]
\Require User key pair $(\mathit{sk}, \mathit{pk})$, last sync timestamp $t_{\mathit{last}}$
\Ensure All pending messages retrieved and decrypted
\State $\mathit{nodes} \leftarrow \mathsf{FindMailbox}(\HH(\mathit{pk}))$
\State $\mathit{msgs} \leftarrow \emptyset$
\For{$n \in \mathit{nodes}$}
    \State $\mathit{batch} \leftarrow \mathsf{Fetch}(n, \HH(\mathit{pk}), t_{\mathit{last}})$
    \State $\mathit{msgs} \leftarrow \mathit{msgs} \cup \mathit{batch}$
\EndFor
\State $\mathit{msgs} \leftarrow \mathsf{Deduplicate}(\mathit{msgs})$ \Comment{By $\mathit{msg\_id}$}
\State $\mathit{msgs} \leftarrow \mathsf{SortByTimestamp}(\mathit{msgs})$
\For{$m \in \mathit{msgs}$}
    \State $p \leftarrow \mathsf{Decrypt}(\mathit{sk}, m.\mathit{ciphertext})$
    \State $\mathsf{Deliver}(p)$
\EndFor
\State $\mathsf{AcknowledgeRetrieval}(\mathit{nodes}, \mathit{msgs})$
\end{algorithmic}
\end{algorithm}

\subsection{Anti-Spam System}

\protname employs a layered anti-spam system that operates without personal identity verification.

\subsubsection{Layer 1: Proof-of-Work}

Every first-contact message must include a proof-of-work token:

\begin{definition}[Contact PoW]
A valid proof-of-work for sender $s$ contacting recipient $r$ is a nonce $w$ such that:
\begin{equation}
    \HH(s.\mathit{pk} \| r.\mathit{pk} \| w \| \lfloor t / T_w \rfloor) < D
\end{equation}
where $D$ is the difficulty target and $T_w$ is the time window (default: 1 hour). The difficulty $D$ is calibrated to require approximately 5 seconds of computation on a modern mobile device.
\end{definition}

\subsubsection{Layer 2: Social Trust Amplification}

Users who are vouched for by existing contacts receive reduced PoW requirements:

\begin{equation}
    D_{\mathit{effective}}(s, r) = D \cdot 2^{\min(\mathit{trust}(s, r), 10)}
\end{equation}

where $\mathit{trust}(s, r)$ is the shortest path length in the social trust graph between $s$ and $r$, and the exponent is capped at 10 to prevent complete PoW elimination.

\subsubsection{Layer 3: Rate Limiting}

Each relay node enforces per-sender rate limits based on accumulated PoW credits:

\begin{algorithm}[H]
\caption{Rate Limit Check}
\label{alg:rate-limit}
\begin{algorithmic}[1]
\Require Sender $s$, message $m$, relay node $n$
\Ensure Accept or reject message
\State $\mathit{credits} \leftarrow n.\mathsf{GetCredits}(s.\mathit{pk})$
\State $\mathit{cost} \leftarrow \mathsf{MessageCost}(|m|)$ \Comment{Proportional to size}
\If{$\mathit{credits} \geq \mathit{cost}$}
    \State $n.\mathsf{DeductCredits}(s.\mathit{pk}, \mathit{cost})$
    \State \Return \textsc{Accept}
\Else
    \State \Return \textsc{Reject}
\EndIf
\end{algorithmic}
\end{algorithm}

\subsection{Media Sharing}

Large media files (images, voice messages, videos) are handled through a content-addressed storage layer integrated with the DSF network:

\begin{enumerate}
    \item The sender encrypts the media file with a random symmetric key $k_m$.
    \item The encrypted media is split into fixed-size chunks (256 KB) and distributed across multiple relay nodes using erasure coding (Reed-Solomon with $n = 2k$ redundancy).
    \item The message body contains the decryption key $k_m$, the content hash, and the chunk manifest.
    \item The recipient retrieves chunks from available nodes and reconstructs the media.
\end{enumerate}

This design ensures that no single relay node possesses the complete encrypted media, and the decryption key is protected by the end-to-end encryption of the message itself.

\section{Security Analysis}

\subsection{Forward Secrecy}

\begin{theorem}[Forward Secrecy]
\label{thm:forward-secrecy}
If the DDH assumption holds in $\GG$, then \protname provides forward secrecy: compromise of a user's long-term identity key does not reveal the contents of previously exchanged messages.
\end{theorem}

\begin{proof}
Forward secrecy follows from the DH ratchet construction. Each ratchet step generates an ephemeral DH key pair $(\mathit{sk}_e, \mathit{pk}_e) \leftarrow \mathsf{KeyGen}(1^\secparam)$ and derives the chain key from $\mathsf{DH}(\mathit{sk}_e, \mathit{pk}_r)$ where $\mathit{pk}_r$ is the peer's current ratchet public key. After the ratchet step, $\mathit{sk}_e$ is erased. An adversary who later obtains the long-term key cannot compute $\mathsf{DH}(\mathit{sk}_e, \mathit{pk}_r)$ without $\mathit{sk}_e$, which is equivalent to solving the CDH problem in $\GG$. Under DDH, CDH is hard, completing the proof.
\end{proof}

\subsection{Post-Compromise Security}

\begin{theorem}[Post-Compromise Security]
\label{thm:pcs}
If the DDH assumption holds, \protname provides post-compromise security: after a device compromise, security is restored within at most two DH ratchet steps following the compromise.
\end{theorem}

\begin{proof}
Suppose the adversary compromises the device state at time $t$, obtaining the current ratchet state $(s_t, \mathit{dh}_t)$. After the next DH ratchet step, a fresh key pair $\mathit{dh}_{t+1}$ is generated and the new root key $s_{t+1}$ incorporates $\mathsf{DH}(\mathit{dh}_{t+1}.\mathit{sk}, \mathit{pk}_{\mathit{peer}})$. Since $\mathit{dh}_{t+1}.\mathit{sk}$ is unknown to the adversary (generated after compromise), computing $s_{t+1}$ requires solving CDH. After the peer also performs a ratchet step (generating $\mathit{dh}_{t+2}$ on their side), both directions are secured with keys unknown to the adversary.
\end{proof}

\subsection{Group Security Properties}

\begin{property}[Group Forward Secrecy]
Within an epoch, the symmetric sender key ratchet provides forward secrecy for individual messages. Across epochs, the full re-keying procedure provides forward secrecy against an adversary who compromises a future epoch's group key.
\end{property}

\begin{property}[Group Post-Compromise Security]
Epoch transitions triggered by the time-based policy ($T_e$) ensure that group security recovers from member compromise within at most $T_e$ time, assuming the compromised member does not collude with the adversary beyond the compromise point.
\end{property}

\subsection{Metadata Protection}

\begin{theorem}[Relay Node Metadata Limitation]
A relay node that is not a mailbox node for both sender and recipient learns at most the recipient's mailbox identifier (a hash of their public key) and the encrypted message size. It cannot determine the sender's identity.
\end{theorem}

\begin{proof}
Messages deposited in the DSF network are addressed by $\HH(\mathit{pk}_r)$, which reveals no information about the sender. The message body, including sender identification, is encrypted under the end-to-end encryption layer. The relay node sees only the envelope: $(\mathit{msg\_id}, c, \mathit{rid}, \mathit{timestamp}, \mathit{ttl})$, where $c$ is the ciphertext and $\mathit{rid} = \HH(\mathit{pk}_r)$.
\end{proof}

\section{Persian Web Traffic Camouflage}

A key contribution of \protname is a pluggable transport designed to blend messaging traffic with legitimate Persian web browsing patterns.

\subsection{Traffic Profile Analysis}

We analyzed traffic patterns from 500 consenting users over 30 days to characterize legitimate Persian web traffic. Key findings:

\begin{table}[H]
\centering
\caption{Persian Web Traffic Characteristics}
\label{tab:traffic}
\begin{tabular}{lcc}
\toprule
\textbf{Characteristic} & \textbf{Mean} & \textbf{Std Dev} \\
\midrule
Packet size (bytes) & 847 & 312 \\
Inter-packet delay (ms) & 234 & 187 \\
Session duration (min) & 12.3 & 8.7 \\
Requests per session & 47 & 31 \\
TLS record size (bytes) & 1{,}203 & 456 \\
\bottomrule
\end{tabular}
\end{table}

\subsection{Camouflage Transport Design}

The camouflage transport operates by:

\begin{enumerate}
    \item Establishing a TLS 1.3 connection with an ESNI (Encrypted Server Name Indication) extension to a cooperating front server hosting legitimate Persian content.
    \item Encoding messaging data into HTTP/2 request/response streams that statistically match the traffic profile in Table~\ref{tab:traffic}.
    \item Injecting cover traffic during idle periods to maintain consistent traffic patterns.
    \item Using steganographic encoding in image requests to popular Persian image hosting services as a fallback channel.
\end{enumerate}

\section{Implementation}

\subsection{Cryptographic Primitives}

\begin{table}[H]
\centering
\caption{Cryptographic Primitive Selection}
\label{tab:primitives}
\begin{tabular}{lll}
\toprule
\textbf{Function} & \textbf{Primitive} & \textbf{Rationale} \\
\midrule
Key Agreement & X25519~\citep{bernstein2006curve25519} & Performance, security \\
Signatures & Ed25519~\citep{bernstein2012ed25519} & Compact, fast \\
AEAD & XChaCha20-Poly1305 & Nonce-misuse resistance \\
Hash & BLAKE3 & Speed, parallelism \\
KDF & HKDF-BLAKE3 & Consistency with hash \\
SSS & GF($2^8$) Shamir & Standard, well-analyzed \\
Erasure Coding & Reed-Solomon (16,32) & Mature implementations \\
\bottomrule
\end{tabular}
\end{table}

\subsection{Client Implementation}

The \protname client is implemented as a cross-platform library in Rust with bindings for iOS (Swift), Android (Kotlin), and desktop (Electron/TypeScript). Key implementation details:

\begin{itemize}
    \item \textbf{State Management:} All cryptographic state is stored in an encrypted SQLite database using SQLCipher, with the database key derived from the user's device passcode via Argon2id~\citep{biryukov2016argon2}.
    \item \textbf{Memory Safety:} Sensitive key material is allocated on locked, non-swappable memory pages and securely erased (zeroized) when no longer needed.
    \item \textbf{Binary Size:} The core library compiles to 2.8 MB (ARM64), suitable for distribution in bandwidth-constrained environments.
\end{itemize}

\subsection{Relay Node Implementation}

Relay nodes are implemented in Go for deployment simplicity:

\begin{itemize}
    \item \textbf{Storage:} Messages are stored in an embedded BadgerDB~\citep{badgerdb} key-value store with automatic TTL-based expiration.
    \item \textbf{Networking:} Nodes communicate using QUIC~\citep{langley2017quic} for low-latency multiplexed connections with built-in encryption.
    \item \textbf{DHT:} The Kademlia DHT is implemented using libp2p~\citep{dias2022libp2p} with custom extensions for mailbox routing.
    \item \textbf{Resource Requirements:} A relay node requires approximately 512 MB RAM and 10 GB disk space to serve 10{,}000 mailboxes.
\end{itemize}

\section{Evaluation}

\subsection{Testnet Deployment}

We deployed a testnet of 2{,}400 nodes across 14 countries: United States (480), Canada (360), Germany (240), United Kingdom (240), Sweden (120), Australia (120), Turkey (120), UAE (120), Iran (via cooperating VPN endpoints, 120), France (120), Netherlands (120), Austria (60), Norway (60), and Japan (60).

\subsection{Message Delivery Latency}

\begin{table}[H]
\centering
\caption{Message Delivery Latency (milliseconds)}
\label{tab:latency}
\begin{tabular}{lccccc}
\toprule
\textbf{Scenario} & \textbf{P50} & \textbf{P90} & \textbf{P95} & \textbf{P99} & \textbf{Max} \\
\midrule
Both online, same region & 87 & 142 & 178 & 312 & 890 \\
Both online, cross-region & 340 & 520 & 680 & 1{,}240 & 3{,}100 \\
Recipient offline 1h & 3{,}612 & 5{,}230 & 6{,}100 & 8{,}900 & 12{,}400 \\
Recipient offline 24h & 86{,}450 & 86{,}890 & 87{,}200 & 88{,}100 & 92{,}000 \\
Recipient offline 72h & 259{,}300 & 259{,}800 & 260{,}100 & 261{,}000 & 268{,}000 \\
Via camouflage transport & 1{,}240 & 2{,}100 & 2{,}890 & 4{,}500 & 8{,}200 \\
\bottomrule
\end{tabular}
\end{table}

\subsection{Message Delivery Reliability}

Over a 30-day evaluation period with simulated offline periods:

\begin{table}[H]
\centering
\caption{Message Delivery Reliability}
\label{tab:reliability}
\begin{tabular}{lcc}
\toprule
\textbf{Offline Duration} & \textbf{Delivery Rate} & \textbf{Duplicate Rate} \\
\midrule
0 (always online) & 100.0\% & 0.0\% \\
$\leq$ 1 hour & 100.0\% & 0.1\% \\
$\leq$ 24 hours & 99.97\% & 0.3\% \\
$\leq$ 72 hours & 99.7\% & 0.8\% \\
$\leq$ 7 days & 98.2\% & 1.4\% \\
\bottomrule
\end{tabular}
\end{table}

\subsection{Anti-Spam Effectiveness}

We simulated spam attacks with varying strategies:

\begin{table}[H]
\centering
\caption{Anti-Spam System Performance}
\label{tab:antispam}
\begin{tabular}{lccc}
\toprule
\textbf{Attack Strategy} & \textbf{Messages/hr} & \textbf{Blocked} & \textbf{False Pos.} \\
\midrule
Brute force (no PoW) & 10{,}000 & 100.0\% & 0.0\% \\
PoW-equipped, no trust & 200 & 94.3\% & 0.0\% \\
PoW + fake trust graph & 500 & 99.2\% & 0.1\% \\
Sybil (1{,}000 identities) & 5{,}000 & 98.8\% & 0.2\% \\
Legitimate first contact & 50 & 0.0\% & 0.0\% \\
Legitimate established & 500 & 0.0\% & 0.0\% \\
\bottomrule
\end{tabular}
\end{table}

\subsection{Group Messaging Performance}

\begin{table}[H]
\centering
\caption{Group Epoch Transition Overhead}
\label{tab:group-perf}
\begin{tabular}{lccc}
\toprule
\textbf{Group Size} & \textbf{Rekey Time (ms)} & \textbf{Bandwidth (KB)} & \textbf{Messages/epoch} \\
\midrule
5 & 12 & 2.4 & 847 \\
20 & 45 & 9.6 & 3{,}420 \\
50 & 110 & 24.0 & 8{,}100 \\
100 & 230 & 48.0 & 15{,}600 \\
500 & 1{,}180 & 240.0 & 72{,}000 \\
1{,}000 & 2{,}400 & 480.0 & 134{,}000 \\
\bottomrule
\end{tabular}
\end{table}

\subsection{Comparison with Existing Systems}

\begin{table}[H]
\centering
\caption{Feature Comparison with Existing Messaging Systems}
\label{tab:comparison}
\begin{tabularx}{\textwidth}{lXXXXX}
\toprule
\textbf{Feature} & \textbf{\protname} & \textbf{Signal} & \textbf{Matrix} & \textbf{Briar} & \textbf{Session} \\
\midrule
E2E Encryption & Yes & Yes & Yes (Olm) & Yes & Yes \\
Forward Secrecy & Yes & Yes & Yes & Yes & No \\
Post-Compromise & Yes & Yes & No & Yes & No \\
Decentralized & Yes & No & Partial & Yes & Yes \\
Offline Delivery & Yes (7d) & No & Yes & No & Yes (14d) \\
No Phone/Email & Yes & No & Yes & Yes & Yes \\
Group PCS & Yes & No & No & No & No \\
Traffic Camouflage & Yes & No & No & Tor & No \\
Anti-Spam (no ID) & Yes & No & Partial & Yes & Partial \\
\bottomrule
\end{tabularx}
\end{table}

\section{Deployment Considerations}

\subsection{Relay Node Incentives}

Relay node operators are compensated through the Pars network's ASHA token system. Each successfully delivered message generates a micro-payment to the relay nodes involved in its delivery chain. The payment is proportional to storage duration and bandwidth consumed, with a current rate of approximately 0.001 ASHA per message-day of storage.

\subsection{Client Distribution}

In jurisdictions where app store distribution is restricted, \protname clients are distributed through:

\begin{enumerate}
    \item Direct APK/IPA distribution via the Pars network's decentralized content delivery system.
    \item Domain-fronted download endpoints that appear as legitimate HTTPS traffic to popular CDN providers.
    \item Physical distribution via Bluetooth and local Wi-Fi sharing between trusted contacts.
    \item QR code-based progressive web app (PWA) installation for minimal-footprint deployment.
\end{enumerate}

\subsection{Regulatory Compliance}

\protname is designed to comply with lawful intercept requirements in democratic jurisdictions through transparent architecture: since the system is end-to-end encrypted with no backdoors, compliance consists of documenting that lawful access requires cooperation of the target user's device, consistent with established legal frameworks for encrypted communications~\citep{abelson2015keys}.

\section{Limitations and Future Work}

\begin{enumerate}
    \item \textbf{Global Traffic Analysis:} \protname does not protect against an adversary capable of simultaneously observing all network links. Mixing networks could address this at the cost of increased latency.
    \item \textbf{Large Groups:} Epoch transitions for groups exceeding 1{,}000 members incur significant overhead. TreeKEM-based approaches from MLS could reduce this to $O(\log n)$.
    \item \textbf{Voice/Video:} Real-time communication requires extensions beyond the current store-and-forward model. We are investigating WebRTC integration with the camouflage transport.
    \item \textbf{Key Transparency:} The current design relies on trust-on-first-use (TOFU) for key verification. A key transparency log~\citep{melara2015coniks} would strengthen this.
    \item \textbf{Quantum Resistance:} The X25519 key agreement is not quantum-resistant. Hybrid key encapsulation with ML-KEM (Kyber)~\citep{bos2018crystals} is planned for a future protocol version.
\end{enumerate}

\section{Conclusion}

We have presented \protname, a messaging protocol designed for the specific needs of the Persian diaspora community operating under adversarial network conditions. By combining an extended Double Ratchet protocol with a distributed store-and-forward network, traffic camouflage, and identity-free anti-spam mechanisms, \protname provides strong security guarantees while maintaining practical usability. Our testnet evaluation demonstrates that the protocol achieves acceptable performance for real-world deployment, with sub-second delivery for online users and reliable delivery for users experiencing extended offline periods.

The \protname protocol is open source under the Apache 2.0 license, and we invite the broader security and privacy community to audit and contribute to its development.

\bibliographystyle{plainnat}

\begin{thebibliography}{35}

\bibitem[Anderson and Murdoch(2013)]{anderson2013dimming}
Anderson, C. and Murdoch, S.~J.
\newblock The dimming of the internet: Detecting throttling as a mechanism of censorship in {Iran}.
\newblock \emph{arXiv preprint arXiv:1306.4361}, 2013.

\bibitem[Angel et~al.(2014)]{angel2015obfs4}
Angel, Y., Kadianakis, G., and Murdoch, S.~J.
\newblock obfs4 specification.
\newblock \emph{Tor Project Technical Report}, 2014.

\bibitem[Abelson et~al.(2015)]{abelson2015keys}
Abelson, H., Anderson, R., Bellovin, S.~M., et~al.
\newblock Keys under doormats: Mandating insecurity by requiring government access to all data and communications.
\newblock \emph{Journal of Cybersecurity}, 1(1):69--79, 2015.

\bibitem[Back(2002)]{back2002hashcash}
Back, A.
\newblock Hashcash -- a denial of service counter-measure.
\newblock Technical report, 2002.

\bibitem[BadgerDB(2023)]{badgerdb}
Dgraph Labs.
\newblock BadgerDB: Fast key-value {DB} in {Go}.
\newblock \url{https://github.com/dgraph-io/badger}, 2023.

\bibitem[Barnes et~al.(2023)]{barnes2023mls}
Barnes, R., Beurdouche, B., Robert, R., et~al.
\newblock The {Messaging Layer Security (MLS)} Protocol.
\newblock \emph{RFC 9420}, 2023.

\bibitem[Bernstein(2006)]{bernstein2006curve25519}
Bernstein, D.~J.
\newblock Curve25519: New {Diffie-Hellman} speed records.
\newblock In \emph{PKC 2006}, pages 207--228, 2006.

\bibitem[Bernstein et~al.(2012)]{bernstein2012ed25519}
Bernstein, D.~J., Duif, N., Lange, T., Schwabe, P., and Yang, B.-Y.
\newblock High-speed high-security signatures.
\newblock \emph{Journal of Cryptographic Engineering}, 2(2):77--89, 2012.

\bibitem[Biryukov et~al.(2016)]{biryukov2016argon2}
Biryukov, A., Dinu, D., and Khovratovich, D.
\newblock Argon2: New generation of memory-hard functions for password hashing and other applications.
\newblock In \emph{IEEE EuroS\&P}, pages 292--302, 2016.

\bibitem[Bos et~al.(2018)]{bos2018crystals}
Bos, J., Ducas, L., Kiltz, E., et~al.
\newblock {CRYSTALS} -- {Kyber}: A {CCA}-secure module-lattice-based {KEM}.
\newblock In \emph{IEEE EuroS\&P}, pages 353--367, 2018.

\bibitem[Bozorgmehr(2007)]{bozorgmehr2007iranian}
Bozorgmehr, M.
\newblock {Iranian} immigration and the changing landscape of {American} diversity.
\newblock In \emph{Iranians in Texas}, Univ. of Texas Press, 2007.

\bibitem[Briar Project(2018)]{briar2018}
Briar Project.
\newblock Briar: Secure messaging, anywhere.
\newblock \url{https://briarproject.org}, 2018.

\bibitem[Chase and Perrin(2020)]{chase2020signal}
Chase, M. and Perrin, T.
\newblock The {Signal} Private Group System.
\newblock Technical Report, Signal Foundation, 2020.

\bibitem[Deibert et~al.(2008)]{deibert2008access}
Deibert, R.~J., Palfrey, J.~G., Rohozinski, R., and Zittrain, J.
\newblock \emph{Access Denied: The Practice and Policy of Global Internet Filtering}.
\newblock MIT Press, 2008.

\bibitem[Dias et~al.(2022)]{dias2022libp2p}
Dias, D., Benet, J., and Protocol Labs.
\newblock libp2p: A modular network stack.
\newblock \url{https://libp2p.io}, 2022.

\bibitem[Dingledine et~al.(2004)]{dingledine2004tor}
Dingledine, R., Mathewson, N., and Syverson, P.
\newblock Tor: The second-generation onion router.
\newblock In \emph{USENIX Security Symposium}, 2004.

\bibitem[EFF(2017)]{eff2017border}
Electronic Frontier Foundation.
\newblock Digital privacy at the {U.S.} border.
\newblock \url{https://www.eff.org/wp/digital-privacy-us-border}, 2017.

\bibitem[Ermoshina et~al.(2017)]{ermoshina2017can}
Ermoshina, K., Musiani, F., and Halpin, H.
\newblock End-to-end encrypted messaging protocols: An overview.
\newblock In \emph{INSCI}, pages 244--254, 2017.

\bibitem[Fifield et~al.(2015)]{fifield2015blocking}
Fifield, D., Lan, C., Hynes, R., Wegmann, P., and Paxson, V.
\newblock Blocking-resistant communication through domain fronting.
\newblock In \emph{PETS}, pages 46--64, 2015.

\bibitem[Fifield(2015)]{fifield2015meek}
Fifield, D.
\newblock meek pluggable transport.
\newblock Technical report, Tor Project, 2015.

\bibitem[Jacob and Eikenberg(2019)]{jacob2019matrix}
Jacob, F. and Eikenberg, R.
\newblock Analysis of the {Matrix} event graph replicated data type.
\newblock \emph{arXiv preprint arXiv:1907.12539}, 2019.

\bibitem[Karlin et~al.(2011)]{karlin2011decoy}
Karlin, J., Ellard, D., Jackson, A.~W., et~al.
\newblock Decoy routing: Toward unblockable internet communication.
\newblock In \emph{USENIX FOCI}, 2011.

\bibitem[Langley et~al.(2017)]{langley2017quic}
Langley, A., Riddoch, A., Wilk, A., et~al.
\newblock The {QUIC} transport protocol: Design and internet-scale deployment.
\newblock In \emph{ACM SIGCOMM}, pages 183--196, 2017.

\bibitem[Marczak et~al.(2014)]{marczak2014hacking}
Marczak, B., Guarnieri, C., Marquis-Boire, M., and Scott-Railton, J.
\newblock Hacking {Team} and the targeting of {Ethiopian} journalists.
\newblock \emph{Citizen Lab Research Report}, 2014.

\bibitem[Marlinspike and Perrin(2016)]{marlinspike2016signal}
Marlinspike, M. and Perrin, T.
\newblock The {Signal} Protocol.
\newblock Technical report, Signal Foundation, 2016.

\bibitem[Maymounkov and Mazi\`{e}res(2002)]{maymounkov2002kademlia}
Maymounkov, P. and Mazi\`{e}res, D.
\newblock Kademlia: A peer-to-peer information system based on the {XOR} metric.
\newblock In \emph{IPTPS}, pages 53--65, 2002.

\bibitem[Melara et~al.(2015)]{melara2015coniks}
Melara, M.~S., Blankstein, A., Bonneau, J., Felten, E.~W., and Freedman, M.~J.
\newblock {CONIKS}: Bringing key transparency to end users.
\newblock In \emph{USENIX Security}, pages 383--398, 2015.

\bibitem[Mislove et~al.(2008)]{mislove2008ostra}
Mislove, A., Post, A., Gummadi, K.~P., and Druschel, P.
\newblock Ostra: Leveraging trust to thwart unwanted communication.
\newblock In \emph{NSDI}, pages 15--30, 2008.

\bibitem[Perrin and Marlinspike(2016)]{perrin2016double}
Perrin, T. and Marlinspike, M.
\newblock The Double Ratchet Algorithm.
\newblock Technical Report, Signal Foundation, 2016.

\bibitem[Shamir(1979)]{shamir1979share}
Shamir, A.
\newblock How to share a secret.
\newblock \emph{Communications of the ACM}, 22(11):612--613, 1979.

\bibitem[Winter and Lindskog(2012)]{winter2012great}
Winter, P. and Lindskog, S.
\newblock How the {Great Firewall of China} is blocking {Tor}.
\newblock In \emph{USENIX FOCI}, 2012.

\bibitem[Wustrow et~al.(2011)]{wustrow2011telex}
Wustrow, E., Wolchok, S., Goldberg, I., and Halderman, J.~A.
\newblock Telex: Anticensorship in the network infrastructure.
\newblock In \emph{USENIX Security}, pages 459--474, 2011.

\end{thebibliography}

\appendix

\section{Protocol Message Formats}

\subsection{Envelope Format}

\begin{lstlisting}[basicstyle=\small\ttfamily,frame=single,caption={DSF Envelope Binary Format}]
Envelope {
    version:    uint8       // Protocol version (1)
    msg_id:     [32]byte    // BLAKE3 hash
    recipient:  [32]byte    // H(pk_recipient)
    timestamp:  uint64      // Unix timestamp (ms)
    ttl:        uint32      // Time-to-live (seconds)
    flags:      uint16      // Delivery flags
    payload:    []byte      // Encrypted message
    pow_nonce:  [8]byte     // Anti-spam PoW
    pow_hash:   [32]byte    // PoW verification
}
\end{lstlisting}

\subsection{Message Format}

\begin{lstlisting}[basicstyle=\small\ttfamily,frame=single,caption={Encrypted Message Format}]
Message {
    header:     []byte      // Encrypted DR header
    ciphertext: []byte      // AEAD encrypted body
    media_refs: []MediaRef  // Optional media references
}

MediaRef {
    content_hash: [32]byte  // BLAKE3 of plaintext
    enc_key:      [32]byte  // Symmetric media key
    chunk_count:  uint16    // Number of chunks
    total_size:   uint64    // Original file size
    mime_type:    string    // MIME type
    manifest:     []ChunkLoc // Chunk locations
}
\end{lstlisting}

\section{Trust Graph Construction}

The social trust graph is constructed through mutual key verification ceremonies. When two users verify each other's keys (via QR code scanning, safety number comparison, or in-person verification), a bidirectional trust edge is created. Trust edges are stored locally and propagated through the network using a gossip protocol with the following properties:

\begin{itemize}
    \item Trust edges are signed by both parties using their identity keys.
    \item Trust propagation is limited to 3 hops to bound graph exploration.
    \item Trust edge creation is rate-limited to 10 new edges per day per user.
    \item Revocation of trust edges takes effect within 1 hour across the network.
\end{itemize}

\end{document}
